% Composable, domain-neutral content components for the resume template.
%
% Load this file from the preamble after theme.tex. The public API deliberately
% models visual composition (slots, cards, grids, and lists), not resume fields.
% Existing content can keep using the compatibility commands at the end.

\RequirePackage{xurl}
\RequirePackage{needspace}
\RequirePackage{varwidth}
\tcbuselibrary{breakable,raster,skins}

% Safe layout defaults. A layout engine may define any of these before loading
% this file to change density without changing content.
\providecommand{\ResumeHeaderKicker}{PROFILE}
\providecommand{\ResumeHeaderKeepSpace}{24mm}
\providecommand{\ResumeHeaderAfterSkip}{1.50mm}
\providecommand{\ResumeHeaderColumnGap}{4.2mm}
\providecommand{\ResumeContactColumnGap}{2.6mm}
\providecommand{\ResumeContactRowGap}{1.15mm}
\providecommand{\ResumeSectionNeedspace}{14mm}
\providecommand{\ResumeEntryNeedspace}{11mm}
\providecommand{\ResumeBandNeedspace}{32mm}
\providecommand{\ResumeSectionBefore}{1.20mm}
\providecommand{\ResumeSectionAfter}{0.50mm}
\providecommand{\ResumeGridGap}{4.2mm}
\providecommand{\ResumeGridGapTwo}{4.65mm}
\providecommand{\ResumeGridGapThree}{4.15mm}
\providecommand{\ResumeGridRowGap}{1.55mm}
\providecommand{\ResumeCardBefore}{0pt}
\providecommand{\ResumeCardAfter}{0pt}
\providecommand{\ResumeCardPaddingX}{1.8mm}
\providecommand{\ResumeCardPaddingY}{1.15mm}
\providecommand{\ResumeEntryBefore}{1.10mm}
\providecommand{\ResumeEntryAfter}{0.60mm}
\providecommand{\ResumeProjectNeedspace}{11mm}
\providecommand{\ResumeProjectBefore}{0.80mm}
\providecommand{\ResumeProjectAfter}{0.35mm}
\providecommand{\ResumeEntryColumnGap}{2.4mm}
\providecommand{\ResumeMetaColumnGap}{2.2mm}
\providecommand{\ResumeListLeftMargin}{4.45mm}
\providecommand{\ResumeListItemSep}{0.50mm}
\providecommand{\ResumeListTopSep}{0.52mm}
\providecommand{\ResumeChipMaxWidth}{42mm}
\providecommand{\ResumeHeaderNameSize}{23.2}
\providecommand{\ResumeHeaderNameLeading}{24.4}
\providecommand{\ResumeHeadlineSize}{9.7}
\providecommand{\ResumeHeadlineLeading}{11.1}
\providecommand{\ResumeSectionTitleSize}{11.9}
\providecommand{\ResumeSectionTitleLeading}{12.6}
\providecommand{\ResumeEntryTitleSize}{10.2}
\providecommand{\ResumeEntryTitleLeading}{11.2}
\providecommand{\ResumeDetailSize}{8.9}
\providecommand{\ResumeDetailLeading}{10.2}
\providecommand{\ResumeMetaSize}{8.35}
\providecommand{\ResumeMetaLeading}{9.5}

% Column weights must add up to 2 because tabularx sees two X columns.
\providecommand{\ResumeHeaderLeftWeight}{0.88}
\providecommand{\ResumeHeaderRightWeight}{1.12}
\providecommand{\ResumeEntryLeftWeight}{1.36}
\providecommand{\ResumeEntryRightWeight}{0.64}
\providecommand{\ResumeMetaLeftWeight}{1.24}
\providecommand{\ResumeMetaRightWeight}{0.76}

\providecommand{\ResumeNeedspace}[1]{%
  \ifinner\else\Needspace*{#1}\fi}
\providecommand{\ResumeNeedSection}{%
  \ResumeNeedspace{\ResumeSectionNeedspace}}
\providecommand{\ResumeNeedSectionBand}{%
  \ResumeNeedspace{%
    \dimexpr\ResumeSectionNeedspace+\ResumeBandNeedspace\relax}}
\providecommand{\ResumeNeedEntry}{%
  \ResumeNeedspace{\ResumeEntryNeedspace}}
\providecommand{\ResumeNeedBand}{%
  \ResumeNeedspace{\ResumeBandNeedspace}}
\providecommand{\ResumeNeedProject}{%
  \ResumeNeedspace{\ResumeProjectNeedspace}}

\providecommand{\ResumeRaggedRight}{%
  \raggedright\arraybackslash\hspace{0pt}\setlength{\emergencystretch}{2em}}
\providecommand{\ResumeRaggedLeft}{%
  \raggedleft\arraybackslash\hspace{0pt}\setlength{\emergencystretch}{2em}}

% Inline emphasis helpers remain useful in arbitrary sections.
\providecommand{\lead}[1]{{\color{Ink}\bfseries #1}}
\providecommand{\metric}[1]{{\color{AccentStrong}\bfseries #1}}
\providecommand{\plainlink}[2]{\href{#1}{{\color{Link}#2}}}
\providecommand{\headerlink}[2]{\href{#1}{{\color{Link}\bfseries #2}}}
\providecommand{\dash}{\,\textendash\,}
\providecommand{\dotsep}{\;\textperiodcentered\;}
\providecommand{\thindot}{\,\textperiodcentered\,}
\providecommand{\arrow}{\;$\rightarrow$\;}
\providecommand{\extlink}{\,\raisebox{-0.1ex}{$\scriptstyle\nearrow$}}

% A reusable, zero-width semantic page marker. \label emits only a deferred
% AUX write, so the recorded page is the page that actually ships while the
% rendered document receives no glyph, box, glue, or hyperlink target.
\ProvideDocumentCommand{\resumepagemarker}{m}{%
  \label{resume-page-marker:#1}}

% Typography atoms. These never create a box, so they can wrap in any slot.
\ProvideDocumentCommand{\resumeentrytitle}{m}{%
  {\color{AccentStrong}%
   \fontsize{\ResumeEntryTitleSize}{\ResumeEntryTitleLeading}\selectfont
   \bfseries #1}}
\ProvideDocumentCommand{\resumeentrydetail}{m}{%
  {\color{Muted}%
   \fontsize{\ResumeDetailSize}{\ResumeDetailLeading}\selectfont #1}}
\ProvideDocumentCommand{\resumedate}{m}{%
  {\color{Ink}%
   \fontsize{\ResumeMetaSize}{\ResumeMetaLeading}\selectfont\bfseries #1}}
\ProvideDocumentCommand{\resumemuted}{m}{%
  {\color{Muted}%
   \fontsize{\ResumeMetaSize}{\ResumeMetaLeading}\selectfont #1}}

% Chips use varwidth so a long label wraps inside the current column instead of
% widening the page. Short labels retain their natural width.
\ProvideDocumentCommand{\resumechip}{m m +m}{%
  \tcbox[
    on line,
    nobeforeafter,
    colback=#1,
    colframe=#2!18!Paper,
    boxrule=0.28pt,
    arc=0.7mm,
    left=1.0mm,
    right=1.0mm,
    top=0.30mm,
    bottom=0.30mm,
    boxsep=0pt
  ]{%
    \begin{varwidth}{\ResumeChipMaxWidth}
      \raggedright
      \color{#2}\fontsize{7.45}{8.2}\selectfont\bfseries #3%
    \end{varwidth}}}
\ProvideDocumentCommand{\resumetagchip}{+m}{%
  \resumechip{AccentSoft}{AccentStrong}{#1}}
\ProvideDocumentCommand{\resumerolechip}{+m}{%
  \resumechip{SecondarySoft}{Secondary}{#1}}
\ProvideDocumentCommand{\resumemedalchip}{m +m}{%
  \resumechip{#1!11!Paper}{#1}{#2}}
\ProvideDocumentCommand{\resumelinkchip}{m +m}{%
  \href{#1}{\resumechip{AccentSoft}{Link}{\nolinkurl{#2}\extlink}}}
\providecommand{\chip}[3]{\resumechip{#1}{#2}{#3}}
\providecommand{\tagchip}[1]{\resumetagchip{#1}}
\providecommand{\rolechip}[1]{\resumerolechip{#1}}
\providecommand{\medalchip}[2]{\resumemedalchip{#1}{#2}}
\providecommand{\linkchip}[2]{\resumelinkchip{#1}{#2}}

% A mark is intentionally narrow, but its content still wraps instead of
% overflowing. Prefer a short glyph; a long label simply grows vertically.
\ProvideDocumentCommand{\resumemark}{O{Accent} +m}{%
  \tcbox[
    on line,
    nobeforeafter,
    colback=#1!9!Paper,
    colframe=#1!18!Paper,
    boxrule=0.32pt,
    arc=0.9mm,
    left=0.35mm,
    right=0.35mm,
    top=0.55mm,
    bottom=0.50mm,
    boxsep=0pt
  ]{%
    \parbox{4.4mm}{%
      \centering\color{#1}\fontsize{7.0}{7.5}\selectfont\bfseries #2}}}
\providecommand{\brandmark}[2][Accent]{\resumemark[#1]{#2}}

% Generic 1/2/3-column raster. It accepts any number of resumegriditem boxes;
% tcolorbox starts a new row automatically when the selected column count is
% reached. Override raster equal height=none for deliberately uneven content.
\ProvideDocumentEnvironment{resumegrid}{O{2} O{}}{%
  \def\ResumeResolvedGridGap{\ResumeGridGap}%
  \ifnum#1=2\relax
    \def\ResumeResolvedGridGap{\ResumeGridGapTwo}%
  \fi
  \ifnum#1=3\relax
    \def\ResumeResolvedGridGap{\ResumeGridGapThree}%
  \fi
  \begin{tcbraster}[
    raster columns=#1,
    raster column skip=\ResumeResolvedGridGap,
    raster row skip=\ResumeGridRowGap,
    raster equal height=rows,
    raster valign=top,
    raster before skip=0pt,
    raster after skip=0pt,
    #2
  ]%
}{%
  \end{tcbraster}%
}
\ProvideDocumentCommand{\resumegriditem}{+m}{%
  \begin{tcolorbox}[
    enhanced,
    blankest,
    boxsep=0pt,
    left=0pt,
    right=0pt,
    top=0pt,
    bottom=0pt,
    before skip=0pt,
    after skip=0pt
  ]
    \ResumeRaggedRight #1%
  \end{tcolorbox}}

% Variable-length contact grid and header. The body is a sequence of
% \resumecontact items; the optional column count may be 1, 2, or 3.
\ProvideDocumentEnvironment{resumecontactgrid}{O{2}}{%
  \begin{tcbraster}[
    raster columns=#1,
    raster column skip=\ResumeContactColumnGap,
    raster row skip=\ResumeContactRowGap,
    raster equal height=rows,
    raster valign=top,
    raster before skip=0pt,
    raster after skip=0pt
  ]%
}{%
  \end{tcbraster}%
}
\providecommand{\contactlabel}[1]{%
  {\color{Subtle}\fontsize{6.6}{7.2}\selectfont\bfseries
   \addfontfeature{LetterSpace=18}\MakeUppercase{#1}}}
\providecommand{\contactvalue}[1]{%
  {\fontsize{8.7}{10.4}\selectfont #1}}
\ProvideDocumentCommand{\resumecontact}{m +m}{%
  \begin{tcolorbox}[
    enhanced,
    blankest,
    boxsep=0pt,
    left=0pt,
    right=0pt,
    top=0pt,
    bottom=0pt,
    before skip=0pt,
    after skip=0pt
  ]
    \raggedleft\hspace{0pt}%
    \contactlabel{#1}\hspace{1.25mm}\contactvalue{#2}%
  \end{tcolorbox}}

\ProvideDocumentCommand{\resumeprofileheader}{O{2} m m O{\ResumeHeaderKicker} +m}{%
  \Needspace{\ResumeHeaderKeepSpace}%
  \begin{minipage}[b]{0.52\linewidth}
    {\fontsize{\ResumeHeaderNameSize}{\ResumeHeaderNameLeading}\selectfont
     \bfseries #2}\par
    \vspace{\ResumeHeaderNameGap}
    {\color{Muted}%
     \fontsize{\ResumeHeadlineSize}{\ResumeHeadlineLeading}\selectfont #3}
  \end{minipage}\hfill
  \begin{minipage}[b]{0.46\linewidth}
    \begin{resumecontactgrid}[#1]
      #5%
    \end{resumecontactgrid}
  \end{minipage}\par
  \vspace{\ResumeHeaderAfterSkip}%
}
\ProvideDocumentEnvironment{resumeheaderblock}{O{2} m m O{\ResumeHeaderKicker} +b}{%
  \resumeprofileheader[#1]{#2}{#3}[#4]{#5}%
}{}

% Arbitrary section title. The optional subtitle may be omitted or may contain
% any short descriptor; both cells wrap through tabularx.
\ProvideDocumentCommand{\resumesection}{m O{}}{%
  \ResumeNeedSection
  \par\addvspace{\ResumeSectionBefore}%
  \begin{tabularx}{\linewidth}{@{}p{1.15mm}@{\hspace{1.45mm}}X@{}}
    {\color{Accent}\rule[-0.15mm]{1.15mm}{2.75mm}} &
    \begin{varwidth}[t]{0.76\linewidth}
      \raggedright
      {\color{AccentStrong}%
       \fontsize{\ResumeSectionTitleSize}{\ResumeSectionTitleLeading}\selectfont
       \bfseries #1}%
      \ifstrempty{#2}{}{%
        \hspace{1.9mm}%
        {\color{Subtle}\fontsize{7.15}{8.2}\selectfont\bfseries
         \addfontfeature{LetterSpace=18}\MakeUppercase{#2}}}%
    \end{varwidth}%
    \hspace{1.4mm}%
    {\color{Rule}%
     \leaders\hrule height 4.0pt depth -3.62pt\hfill\kern0pt}%
  \end{tabularx}%
  \par\nobreak\vspace{\ResumeSectionAfter}\nopagebreak[4]%
}

% Generic flowing card. Its body may contain arbitrary paragraphs, metadata,
% and resumelist environments. The card is breakable when used outside a grid.
\ProvideDocumentEnvironment{resumecard}{O{}}{%
  \begin{tcolorbox}[
    enhanced,
    breakable,
    colback=Surface,
    colframe=Rule,
    boxrule=0.28pt,
    arc=0.9mm,
    left=\ResumeCardPaddingX,
    right=\ResumeCardPaddingX,
    top=\ResumeCardPaddingY,
    bottom=\ResumeCardPaddingY,
    boxsep=0pt,
    before skip=\ResumeCardBefore,
    after skip=\ResumeCardAfter,
    #1
  ]%
}{%
  \end{tcolorbox}%
}

% Generic entry band with arbitrary left and right slots. Both slots are X
% columns and therefore wrap; neither side assumes organization/date fields.
\ProvideDocumentCommand{\resumeentryband}{O{} +m +m}{%
  \ResumeNeedBand
  \begin{tcolorbox}[
    enhanced,
    colback=AccentSurface,
    colframe=Accent,
    boxrule=0pt,
    leftrule=2.05pt,
    arc=0.85mm,
    left=2.45mm,
    right=2.6mm,
    top=1.25mm,
    bottom=1.25mm,
    boxsep=0pt,
    before skip=\ResumeEntryBefore,
    after skip=\ResumeEntryAfter,
    #1
  ]
    \begin{tabularx}{\linewidth}{%
      @{}%
      >{\hsize=\ResumeEntryLeftWeight\hsize\ResumeRaggedRight}X%
      @{\hspace{\ResumeEntryColumnGap}}%
      >{\hsize=\ResumeEntryRightWeight\hsize\ResumeRaggedLeft}X%
      @{}}
      #2 & #3
    \end{tabularx}
  \end{tcolorbox}%
}

% A two-slot metadata line. It is deliberately a paragraph table rather than
% an hbox so long links, dates, labels, and technology summaries can wrap.
\ProvideDocumentCommand{\resumemetarow}{+m +m}{%
  \begin{tabularx}{\linewidth}{%
    @{}%
    >{\hsize=\ResumeMetaLeftWeight\hsize\ResumeRaggedRight}X%
    @{\hspace{\ResumeMetaColumnGap}}%
    >{\hsize=\ResumeMetaRightWeight\hsize\ResumeRaggedLeft}X%
    @{}}
    #1 & #2
  \end{tabularx}%
}

% Compact blue subheading for projects, papers, outcomes, or any other record.
% The right slot is optional in meaning, not in structure, and wraps safely.
\ProvideDocumentCommand{\resumeitemheading}{+m +m}{%
  \ResumeNeedProject
  \vspace{\ResumeProjectBefore}%
  \resumemetarow{%
    {\color{Accent}\rule[-0.28mm]{1.65pt}{3.0mm}}\hspace{1.6mm}%
    \resumeentrytitle{#1}%
  }{%
    #2%
  }%
  \par\vspace{\ResumeProjectAfter}}

\ProvideDocumentEnvironment{resumelist}{O{}}{%
  \begin{itemize}[
    leftmargin=\ResumeListLeftMargin,
    itemsep=\ResumeListItemSep,
    topsep=\ResumeListTopSep,
    parsep=0pt,
    partopsep=0pt,
    #1
  ]%
}{%
  \end{itemize}%
}

% ---------------------------------------------------------------------------
% Compatibility layer
% ---------------------------------------------------------------------------
% Historical names are fallback wrappers only. If an existing entrypoint
% already defines them, \providecommand preserves its visual footprint exactly.
% New content should use the namespaced composable API above.

\providecommand{\rowline}[2]{\resumemetarow{#1}{#2}}
\providecommand{\resumeheader}[6]{%
  \resumeprofileheader[2]{#1}{#2}{%
    \resumecontact{TEL}{#3}%
    \resumecontact{MAIL}{#4}%
    \resumecontact{GITHUB}{#5}%
    \resumecontact{WEB}{#6}%
  }}
\providecommand{\sectiontitle}[2]{\resumesection{#1}[#2]}
\providecommand{\twocol}[2]{%
  \begin{resumegrid}[2]
    \resumegriditem{#1}%
    \resumegriditem{#2}%
  \end{resumegrid}}
\providecommand{\threecol}[3]{%
  \begin{resumegrid}[3]
    \resumegriditem{#1}%
    \resumegriditem{#2}%
    \resumegriditem{#3}%
  \end{resumegrid}}

\ProvideDocumentCommand{\resumelegacyeducation}{m m m m m m}{%
  \begin{tabularx}{\linewidth}{%
    @{}p{5.8mm}@{\hspace{1.3mm}}%
    >{\hsize=1.34\hsize\ResumeRaggedRight}X%
    @{\hspace{1.8mm}}%
    >{\hsize=0.66\hsize\ResumeRaggedLeft}X@{}}
    \resumemark{#1} &
    {\fontsize{9.9}{10.8}\selectfont\bfseries #2} &
    \resumedate{#4}
  \end{tabularx}
  \vspace{0.45mm}%
  \begin{tabularx}{\linewidth}{%
    @{}p{5.8mm}@{\hspace{1.3mm}}%
    >{\hsize=1.18\hsize\ResumeRaggedRight}X%
    @{\hspace{1.8mm}}%
    >{\hsize=0.82\hsize\ResumeRaggedLeft}X@{}}
    {} & \resumeentrydetail{#5} &
    {\resumeentrydetail{#6}\par\vspace{0.25mm}#3}
  \end{tabularx}}
\providecommand{\eduentry}[6]{%
  \resumelegacyeducation{#1}{#2}{#3}{#4}{#5}{#6}}

\ProvideDocumentCommand{\resumelegacyaward}{m m m m}{%
  {\fontsize{8.9}{10.0}\selectfont\bfseries #1}\par
  \vspace{0.55mm}%
  \resumemetarow{\resumemedalchip{#2}{#3}}{\resumemuted{#4}}}
\providecommand{\awardcard}[4]{%
  \resumelegacyaward{#1}{#2}{#3}{#4}}

\ProvideDocumentCommand{\resumelegacyentry}{m m m m m}{%
  \resumeentryband{%
    \resumemark{#1}\hspace{1.45mm}%
    \resumeentrytitle{#2}\hspace{1.45mm}\resumeentrydetail{#3}%
  }{%
    \resumedate{#4}%
    \ifstrempty{#5}{}{\hspace{1.1mm}\resumerolechip{#5}}%
  }}
\providecommand{\cventry}[5]{%
  \resumelegacyentry{#1}{#2}{#3}{#4}{#5}}

\providecommand{\ptick}{%
  {\color{Accent}\rule[-0.28mm]{1.65pt}{3.0mm}}\hspace{1.6mm}}
\ProvideDocumentCommand{\resumeprojectheading}{m m}{%
  \resumeitemheading{#1}{%
    \ifstrempty{#2}{}{\resumerolechip{#2}}%
  }}
\providecommand{\project}[1]{\resumeprojectheading{#1}{}}
\providecommand{\projecttag}[2]{\resumeprojectheading{#1}{#2}}

\ProvideDocumentCommand{\resumelegacypaper}{m m m}{%
  \resumeitemheading{#1}{\resumemuted{#3}}%
  #2}
\providecommand{\paper}[3]{\resumelegacypaper{#1}{#2}{#3}}
